% TOP_PROBLEM: {"id": 7500108, "title": "Increasing Polarized Ramsey Theorem", "category_id": 18, "category": {"id": 18, "name": "logic", "display_name": "Logic", "description": "Problems in mathematical logic, model theory, proof theory, and finite model theory.", "slug": "logic", "order_index": 18, "created_at": "2026-07-31T15:26:25.671Z"}, "status": "partially_solved"}
% FIRST_POSTED: null
% ENTRY_KIND: statement_only
% Imported from: batch20.tex; UnsolvedMath ID 7500108
\documentclass[11pt]{article}
\usepackage[a4paper,margin=25mm]{geometry}
\usepackage{amsmath,amssymb,mathtools}
\usepackage{fontspec,unicode-math}
\setmainfont{Latin Modern Roman}
\setsansfont{Latin Modern Sans}
\setmathfont{Latin Modern Math}
\usepackage{xurl,hyperref}
\hypersetup{colorlinks=true,linkcolor=blue,urlcolor=blue,pdftitle={UnsolvedMath Problem 7500108}}
\newcommand{\sref}[2]{\hyperlink{source:#1:#2}{[S#2]}}
\newcommand{\cataloguescope}[1]{\paragraph{Mathematical question.}#1}
\begin{document}
% UnsolvedMath ID 7500108; source catalogue code AMR-074-0108
\setcounter{section}{7500107}
\section{Increasing Polarized Ramsey Theorem}\label{7500108}
{\small\sffamily UnsolvedMath ID 7500108 · Logic}
\subsection{Definitions and mathematical statement}

\paragraph{Shared notation.}$\mathbb N=\{1,2,\ldots\}$, and $\mathbb Z,\mathbb Q,\mathbb R,\mathbb C$ denote the integers, rationals, reals and complex numbers. Any other conventions are specified locally.


Work in second-order arithmetic over $\mathsf{RCA}_0$, the base theory with $\Delta^0_1$ comprehension and $\Sigma^0_1$ induction. The strength of a principle means which subsystems prove it and which subsystems it implies over that base; equivalence requires both implications. A colouring $f:[\mathbb N]^2\to\{0,1\}$ assigns a colour to each unordered pair; write $f(x,y)$ for $x<y$. The principle $\mathsf{RT}^2_2$ asserts an infinite homogeneous set $H$. A colouring is stable if for every $x$ the values $f(x,y)$ are constant for all sufficiently large $y$; $\mathsf{SRT}^2_2$ restricts $\mathsf{RT}^2_2$ to stable colourings. The increasing polarized principle $\mathsf{IPT}^2_2$ asserts infinite $H_1,H_2\subseteq\mathbb N$ such that all $f(x,y)$ with $x\in H_1$, $y\in H_2$ and $x<y$ have one colour.
\paragraph{Statement recorded in the supplied source.}
Is $\mathsf{IPT}^2_2$ equivalent to $\mathsf{RT}^2_2$, to $\mathsf{SRT}^2_2$, to both, or to neither? \sref{7500108}{2}
\subsection{Short English statement}
Does the increasing polarized Ramsey principle have the same strength as the ordinary two-colour pair theorem, its stable version, both, or neither?
\subsection{Sources}
\begin{description}
\item[\hypertarget{source:7500108:1}{S1}] ULAM.AI and the UnsolvedMath Contributors, UnsolvedMath: A Curated Collection of Open Mathematics Problems, record 7500108 (AMR-074-0108). CC BY 4.0. \url{https://huggingface.co/datasets/ulamai/UnsolvedMath}.
\item[\hypertarget{source:7500108:2}{S2}] Antonio Montalbán, Open Questions in Reverse Mathematics, author preprint dated 17 August 2010, Question 8, printed p. 4. \url{https://math.berkeley.edu/~antonio/papers/questionsRM.pdf}.
\end{description}
\label{end:7500108}
\end{document}
